\documentclass{article}

%%%%%%%%%%%%%%% LIBRERIAS %%%%%%%%%%%%%%%%%%%%%
\usepackage[thinc]{esdiff}
\usepackage{amsmath,amssymb}
\usepackage{caption}
\usepackage{titlesec}
\usepackage{titletoc}
\usepackage{graphicx}
\IfFileExists{spanish.ldf}
  {\usepackage[spanish,es-tabla]{babel}}
  {\usepackage[provide=*,spanish]{babel}}
\usepackage{hyperref}
\usepackage{float}
\usepackage[table]{xcolor}
\usepackage{circuitikz}
\usepackage[left=2cm, right=2cm, top=2cm, bottom=2cm, footskip=1cm]{geometry}
\usepackage{subfiles}
\usepackage{subcaption}
\usepackage{booktabs}
\usepackage{grffile}
\usepackage{comment}
\usepackage{pgfplots}
\pgfplotsset{compat=1.18}
\usepackage{pgfplotstable}
\usepgfplotslibrary{groupplots}

\usepackage{array}
\usepackage{longtable}
\usepackage{tabularx}
\usepackage{enumitem}
\usepackage{setspace}
\usepackage{fancyhdr}
\usetikzlibrary{babel,arrows.meta,calc,positioning}

%%%%%%%%%%%%%%%%%%% VARIABLES %%%%%%%%%%%%%%%%%%%%
\newcommand{\Facultad}{Instituto Tecnológico \\de\\ Buenos Aires}
\newcommand{\TPn}{Trabajo Práctico N° 2}
\newcommand{\TPtema}{Determinación de la carga específica del electrón}
\graphicspath{{imagenes/}}

%%%%%%%%%%%%%%%%%% FORMATO TÍTULO Y NUMERACIÓN %%%%%%%%%%%%%%%%%%%
\renewcommand{\thesection}{\arabic{section}.}
\renewcommand{\thesubsection}{\thesection\arabic{subsection}}
\renewcommand{\thesubsubsection}{$\alph{subsubsection}$}
\setcounter{secnumdepth}{3}

\titleformat{\section}{\Large\bfseries}{\thesection}{1em}{}
\titleformat{\subsection}{\large\bfseries}{\thesubsection}{0.5em}{}
\titleformat{\subsubsection}{\large\bfseries}{\thesubsubsection}{0.5em}{}

%%%%%%%%%%%%%%% SOPORTE DEL CUERPO DEL INFORME %%%%%%%%%%%%%%%%%%%%%
\definecolor{azulencabezado}{HTML}{163A5F}
\definecolor{azulclaro}{HTML}{EAF2F8}
\definecolor{grislinea}{HTML}{D9D9D9}
\definecolor{grispendiente}{HTML}{5C5C5C}

\newcommand{\pendiente}[1]{\textcolor{grispendiente}{\textit{[#1]}}}
\newcolumntype{Y}{>{\raggedright\arraybackslash}X}
\newcolumntype{C}[1]{>{\centering\arraybackslash}p{#1}}
\newcolumntype{L}[1]{>{\raggedright\arraybackslash}p{#1}}

\hypersetup{
  hidelinks,
  pdftitle={Informe del experimento de Determinación de la carga específica del electrón},
  pdfsubject={Laboratorio de Física IV},
  pdfkeywords={Carga especifica, laboratorio, Física IV}
}

\makeatletter
\providecommand{\ttl@finishall}{}
\makeatother

%%%%%%%%%%%%%%%%%%% ARCHIVO %%%%%%%%%%%%%%%%%%%%%%%%
\begin{document}
\renewcommand{\tablename}{Tabla}
\hypersetup{pageanchor=false}

%%%CARATULA%%%
\begin{titlepage} %creo portada

        \begin{flushleft}
            \centering
            \includegraphics[width=0.3\textwidth]{Logo_ITBA.png}
        \end{flushleft}

        \centering
            
        {\scshape\LARGE \Facultad \par} 
        \vspace{1cm} 


        {\huge\bfseries \TPn \par}
        \vspace{1.5cm}
        {\Large Laboratorio de Física IV\\ 93.61 \par}
        \vspace{0.5cm}
        {\Large \TPtema \par}
        \vfill 
        {\Large \bfseries Grupo N° 55 \par}
        \vspace{1cm}
        {\large Juan Ignacio Fogolin Lagares \hfill Legajo: 66409  \par}
        {\large Rita Moschini \hfill Legajo: 67026 \par} 
        {\large Agustín Alejandro Montoto \hfill Legajo: 65511\par}
        \vfill
        {\large Fecha de entrega: \today\par}
        \vfill

        \begin{flushleft}
        \renewcommand{\arraystretch}{2.2}
        \begin{tabularx}{\textwidth}{@{}L{8.5cm}Y@{}}
            \toprule
            \textbf{Fecha de aprobación:} & \\
            \midrule
            \textbf{Comentarios y observaciones del docente:} & \\
             & \\
             & \\
            \bottomrule
        \end{tabularx}
        \end{flushleft}

\end{titlepage}

\pagenumbering{gobble}
\tableofcontents 
\newpage

%%%%%%%%%%%%%%%%%%% CUERPO DEL INFORME %%%%%%%%%%%%%%%%%%%%%%%%
\pagenumbering{arabic}
\setcounter{page}{1}
\hypersetup{pageanchor=true}

\setstretch{1.12}
\setlength{\parindent}{0.7cm}
\setlength{\parskip}{0.25em}
\setlength{\arrayrulewidth}{0.45pt}
\setlength{\tabcolsep}{5pt}
\arrayrulecolor{black}
\renewcommand{\arraystretch}{1.28}

\renewcommand{\thesection}{\arabic{section}}
\renewcommand{\thesubsection}{\thesection.\arabic{subsection}}
\setcounter{secnumdepth}{2}

\titleformat{\section}{\Large\bfseries\color{black}}{\thesection}{0.6em}{}
\titleformat{\subsection}{\large\bfseries\color{black}}{\thesubsection}{0.55em}{}
\titlespacing*{\section}{0pt}{1.1em}{0.45em}
\titlespacing*{\subsection}{0pt}{0.9em}{0.35em}

\captionsetup{font={small,it},labelfont=it,justification=centering}
\setlist[itemize]{leftmargin=1.05cm,itemsep=0.18em,topsep=0.25em}
\setlist[enumerate]{leftmargin=1.05cm,itemsep=0.22em,topsep=0.25em}

\pagestyle{fancy}
\fancyhf{}
\fancyhead[L]{\small\itshape Laboratorio de Física IV}
\fancyhead[R]{\small\itshape Instituto Tecnológico de Buenos Aires}
\fancyfoot[C]{\thepage}
\renewcommand{\headrulewidth}{0pt}
\renewcommand{\footrulewidth}{0pt}

\end{document}